\documentclass[11pt,letterpaper]{article}
\usepackage[T1]{fontenc}
\usepackage[utf8]{inputenc}
\usepackage{amsmath,amsthm}
\usepackage{newpxtext,newpxmath}
\usepackage[margin=1in,headheight=16pt]{geometry}
\usepackage{microtype}
\usepackage{mathtools}
\usepackage[dvipsnames]{xcolor}
\usepackage{booktabs,tabularx,array}
\usepackage{enumitem}
\usepackage{fancyhdr}
\usepackage[most]{tcolorbox}
\usepackage{aliascnt}
\usepackage{hyperref}
\usepackage[nameinlink,noabbrev]{cleveref}
\definecolor{InkBlue}{HTML}{173C59}
\definecolor{Teal}{HTML}{176B71}
\definecolor{LightBlue}{HTML}{EFF5F8}
\hypersetup{colorlinks=true,linkcolor=InkBlue,citecolor=Teal,urlcolor=Teal,
 pdftitle={Reciprocity and Matrix Duality for Preorder Polytopes},
 pdfauthor={Research manuscript prepared with ChatGPT},
 pdfsubject={Proofs of three conjectures of Athanasiadis and Chapoton}}
\pagestyle{fancy}
\fancyhf{}
\fancyhead[L]{\small\textsc{Preorder polytopes}}
\fancyhead[R]{\small Reciprocity and matrix duality}
\fancyfoot[C]{\thepage}
\renewcommand{\headrulewidth}{0.35pt}
\setlength{\parskip}{0.35em}
\setlength{\parindent}{1.2em}
\setcounter{tocdepth}{2}
\allowdisplaybreaks[2]
\newtheorem{theorem}{Theorem}[section]
\newaliascnt{proposition}{theorem}
\newtheorem{proposition}[proposition]{Proposition}
\aliascntresetthe{proposition}
\newaliascnt{lemma}{theorem}
\newtheorem{lemma}[lemma]{Lemma}
\aliascntresetthe{lemma}
\newaliascnt{corollary}{theorem}
\newtheorem{corollary}[corollary]{Corollary}
\aliascntresetthe{corollary}
\theoremstyle{definition}
\newaliascnt{definition}{theorem}
\newtheorem{definition}[definition]{Definition}
\aliascntresetthe{definition}
\newaliascnt{example}{theorem}
\newtheorem{example}[example]{Example}
\aliascntresetthe{example}
\newaliascnt{remark}{theorem}
\newtheorem{remark}[remark]{Remark}
\aliascntresetthe{remark}
\newcommand{\N}{\mathbb N}
\newcommand{\Z}{\mathbb Z}
\newcommand{\R}{\mathbb R}
\newcommand{\Q}{\mathbb Q}
\newcommand{\one}{\mathbf 1}
\newcommand{\zero}{\mathbf 0}
\newcommand{\bvec}{\mathbf b}
\newcommand{\pvec}{\mathbf p}
\newcommand{\qvec}{\mathbf q}
\newcommand{\avec}{\mathbf a}
\newcommand{\xvec}{\mathbf x}
\newcommand{\yvec}{\mathbf y}
\newcommand{\cvec}{\mathbf c}
\newcommand{\JJ}{\mathcal J}
\newcommand{\WW}{\mathcal W}
\newcommand{\DD}{\mathcal D}
\newcommand{\QQ}{\mathcal Q}
\newcommand{\EE}{\mathbb E}
\newcommand{\FF}{\mathbb F}
\newcommand{\ds}{\boldsymbol\Delta}
\DeclareMathOperator{\conv}{conv}
\DeclareMathOperator{\Int}{int}
\DeclareMathOperator{\supp}{supp}
\DeclareMathOperator{\des}{des}
\DeclareMathOperator{\Max}{Max}
\DeclareMathOperator{\Vol}{Vol}
\newcommand{\card}[1]{\lvert #1\rvert}
\newcommand{\sumv}[1]{\lvert #1\rvert}
\newtcolorbox{resultbox}[1][]{colback=LightBlue,colframe=InkBlue,
 boxrule=0.6pt,arc=1pt,left=9pt,right=9pt,top=7pt,bottom=7pt,#1}

\begin{document}
\begin{titlepage}
\thispagestyle{empty}
{\color{Teal}\large\textsc{Research in finite order theory}}\par
\vspace{1.05cm}
{\color{InkBlue}\Huge\bfseries Reciprocity and Matrix Duality\par
\vspace{0.15cm}for Preorder Polytopes\par}
\vspace{0.65cm}
{\Large Proofs of three conjectures\par}
\vspace{0.25cm}
{\large Interior translation, colored arrays, and generalized permutohedra\par}
\vspace{1.0cm}
{\large Prepared for Vladimir Reshetnikov\par}
{\normalsize Research manuscript prepared with ChatGPT\par}
{\normalsize September 19, 2026\par}
\vspace{0.9cm}
\begin{abstract}
We investigate conjectures on the polytopes attached to finite preorders by
Athanasiadis and Chapoton. The central observation is that subtracting the
all-ones vector identifies the interior lattice points of their two-parameter
polytope with the lattice points obtained by decreasing both parameters by one.
Ordinary Ehrhart--Macdonald reciprocity along integer rays then proves their
Conjecture~7.2. We establish an $(n+1)$-parameter strengthening and a second,
falling-binomial formula for the lattice enumerator. This proves their
Conjecture~4.4 on the $h^*$-polynomial and the ordinary, $q=1$ specialization of
Conjecture~4.9.

A separate construction replaces the two color classes of the preorder's
bipartite graph by independently prescribed numbers of copies. Postnikov's
lattice-point formula yields a heterogeneous array duality, which specializes
to the matrix Ehrhart--zeta duality of Conjecture~4.2. Detailed proofs, examples,
sign and boundary checks, and an exact-arithmetic verification program are
provided. The program checks every labeled preorder on one through four
elements. No conjectured shellability or real-rootedness is used.
\end{abstract}
\vfill
\begin{resultbox}
\textbf{Status.} The results are proved in this manuscript relative to the two
explicitly stated classical theorems. The work has not been externally reviewed
or formally verified. The selected statements are conjectures in the public
version of arXiv:2605.26916 examined on September 19, 2026; the literature search
does not establish priority against unpublished work.
\end{resultbox}
\end{titlepage}

\begingroup
\small
\setlength{\parskip}{0pt}
\tableofcontents
\endgroup
\newpage

\section{The selected problem and the outcome}\label{sec:overview}

A finite preorder is a reflexive, transitive relation. Antisymmetry is not
required, so equivalence classes may contain several elements. The operation
of reversing this relation is a natural order-theoretic duality. The question
studied here is how that duality interacts with integer vectors constrained by
order ideals.

The starting point is Conjecture~7.2 of Athanasiadis and Chapoton
\cite{AC26}, in the first public version dated May~26, 2026. For a preorder
$\tau$ on an $n$-element set $E$, write
\begin{equation}\label{eq:uniform-def}
 \QQ_\tau(r,s)=
 \left\{\xvec\in\R_{\geq0}^{E}:
       x(I)\leq r\card I+s\text{ for every nonempty order ideal }I\right\},
 \qquad x(I):=\sum_{e\in I}x_e.
\end{equation}
For $r,s\in\N$, its lattice-point count is a polynomial
$\EE_\tau(r,s)$ of total degree $n$. Their conjecture asks whether
\begin{equation}\label{eq:target}
 \boxed{\EE_\tau(-u,-v)=(-1)^n\EE_\tau(u-1,v-1).}
\end{equation}
Here and throughout, $\N=\{0,1,2,\ldots\}$. Negative arguments refer to
\emph{polynomial evaluation}, not to negative dilations interpreted as feasible
sets.

\subsection{What is proved}
The main results have two independent mechanisms. Interior translation proves
\eqref{eq:target} and its fully weighted extension. A bipartite graph with
independently replicated vertices proves the matrix duality. A binomial
specialization of the first result yields the word formula.

\begin{center}
\begin{tabularx}{\textwidth}{@{}p{0.21\textwidth}X p{0.22\textwidth}@{}}
\toprule
Statement in \cite{AC26} & Result here & Location \\
\midrule
Conjecture 7.2 & Double Ehrhart reciprocity, with a stronger multivariate identity & \Cref{sec:reciprocity} \\
Conjecture 4.2 & Matrix Ehrhart--zeta duality, with heterogeneous copy counts & \Cref{sec:matrix} \\
Conjecture 4.4 & Descent interpretation of the $h^*$-polynomial & \Cref{sec:words} \\
Equation (22) in Conjecture 4.9 & Ordinary zeta evaluation at $-1$ & \Cref{cor:zeta} \\
\bottomrule
\end{tabularx}
\end{center}

The full $q$-identity in Conjecture~4.9 is \emph{not} proved here. Nor do we
claim to settle the paper's questions about shellability, real-rootedness,
$\gamma$-positivity, or support enumerators. Three conjectures are proved, but
one of the three is a consequence of another; these are not three unrelated
breakthroughs.

\subsection{The shortest proof of the original target}
For positive integers $r,s$, the map
\[
 \xvec\longmapsto\xvec-\one
\]
is a bijection
\begin{equation}\label{eq:short-bijection}
 \Int\QQ_\tau(r,s)\cap\Z^E
 \ \longleftrightarrow\
 \QQ_\tau(r-1,s-1)\cap\Z^E.
\end{equation}
Indeed, integrality changes each strict upper inequality into
$x(I)\leq r\card I+s-1$, and subtracting one in each coordinate removes
$\card I$ from its left-hand side. Coordinate positivity becomes
nonnegativity. Apply ordinary Ehrhart reciprocity to the \emph{fixed} lattice
polytope $\QQ_\tau(r,s)$, using
$t\QQ_\tau(r,s)=\QQ_\tau(tr,ts)$. This proves \eqref{eq:target} on all
positive integer pairs, hence as a polynomial identity. The rest of the
article verifies the hypotheses, extends the argument, and develops its
consequences.

\subsection{Provenance and the scope of the novelty claim}
The version history and the statements of Conjectures~4.2, 4.4, and~7.2 in
\cite{AC26} were checked directly, including the rendered PDF pages.
Focused public searches did not locate a resolution. This supports selecting
them as research targets; it does not certify the absence of private or
unindexed proofs. The underlying counting and reciprocity tools are classical.
The contribution of this manuscript is the specialization, the explicit
translation and cloning constructions, and the resulting proofs and extensions.

\section{Notation and the two classical inputs}\label{sec:tools}

\subsection{Preorders, ideals, and finite downsets}
Fix a nonempty finite set $E$, with $n=\card E$, and a preorder
$\leq_\tau$. Its dual $\tau^*$ is defined by
$e\leq_{\tau^*}f$ if and only if $f\leq_\tau e$. An order ideal is a subset
closed downwards. An order filter is closed upwards. Write $\JJ(\tau)$ for
the family of all order ideals, including the empty ideal. In displayed
polytope inequalities we use only the nonempty ideals.

For $S\subseteq E$, let
\[
 D_\tau(S)=\{e:e\leq_\tau f\text{ for some }f\in S\},\qquad
 U_\tau(S)=\{f:e\leq_\tau f\text{ for some }e\in S\}.
\]
Abbreviate $U_\tau(\{e\})$ to $U_\tau(e)$. Reflexivity and transitivity ensure
that these are respectively the ideal and filter generated by $S$.
For any vector $\bvec$, let $b(S)=\sum_{e\in S}b_e$ and
$\sumv{\bvec}=b(E)$.

Define
\begin{equation}\label{eq:basic-P}
 \QQ_\tau:=\QQ_\tau(1,0),\qquad
 P_\tau:=\QQ_\tau\cap\Z^E,
 \qquad B_\tau:=\{\avec\in P_\tau:\sumv{\avec}=n\}.
\end{equation}
The order on $P_\tau$ is coordinatewise comparison. It is a finite downset
in $\N^E$. In particular, if $\avec\in P_\tau$, then the interval
$[\zero,\avec]$ is the product of chains
$\prod_{e\in E}\{0,1,\ldots,a_e\}$.

\begin{lemma}\label{lem:maximal}
The maximal elements of $P_\tau$ are precisely $B_\tau$.
\end{lemma}
\begin{proof}
An element of coordinate sum $n$ cannot be increased because the ideal $E$
provides the bound $x(E)\leq n$. Conversely, suppose that $\avec\in P_\tau$
has $\sumv{\avec}<n$ and cannot be increased in any coordinate. For each
$e\in E$, failure of $\avec+\mathbf e_e$ to be feasible gives an ideal
$I_e$ containing $e$ with $a(I_e)=\card{I_e}$. Such an ideal is called tight.

The union of two tight ideals is tight. To see this, use additivity to write
\[
 a(I\cup J)+a(I\cap J)=a(I)+a(J)=\card I+\card J
 =\card{I\cup J}+\card{I\cap J}.
\]
The two summands on the left are each at most the corresponding cardinalities
on the right. Equality of the sums forces equality for both summands. This
also covers the empty intersection. Repeating the union operation shows that
$\bigcup_{e\in E}I_e=E$ is tight, contradicting $\sumv{\avec}<n$.
Thus some coordinate can be increased, as required.
\end{proof}

\subsection{Ehrhart--Macdonald reciprocity}
If $K\subset\R^d$ is a full-dimensional lattice polytope, its Ehrhart
polynomial $L_K(t)$ satisfies
\begin{equation}\label{eq:EM}
 L_K(t)=\#(tK\cap\Z^d)\quad(t\in\N),\qquad
 L_K(-t)=(-1)^d\#(\Int(tK)\cap\Z^d)\quad(t\geq1).
\end{equation}
We use this classical theorem in the form developed in Stanley's treatment
of Ehrhart--Macdonald reciprocity \cite[Section~6]{Stanley74}. We do not need a
separate multivariate reciprocity theorem: restricting a polynomial to each
positive integer ray is sufficient.

For a full-dimensional $n$-dimensional lattice polytope $K$, define its
$h^*$-polynomial by
\begin{equation}\label{eq:hstar-def}
 \sum_{m\geq0}L_K(m)t^m=\frac{h^*(K,t)}{(1-t)^{n+1}}.
\end{equation}
The finite-difference relation implicit in this definition will be enough
for the computations below.

\subsection{Postnikov's lattice-point formula}
This is the other substantive imported theorem. Let $R$ be a finite set of
size $d+1$, and let $I_0,I_1,\ldots,I_m$ be nonempty subsets of $R$, with
$I_0=R$. Associate to them a bipartite graph with one left vertex for each
$I_i$, right vertex set $R$, and adjacency given by membership. A vector
$\alpha=(\alpha_0,\ldots,\alpha_m)\in\N^{m+1}$ is \emph{draconian} when
\begin{equation}\label{eq:draconian}
 \sum_{i=0}^m\alpha_i=d,\qquad
 \sum_{i\in A}\alpha_i\leq\card{\bigcup_{i\in A}I_i}-1
 \quad\text{for every nonempty }A\subseteq\{0,\ldots,m\}.
\end{equation}
Here $\Delta_I=\conv\{\mathbf e_j:j\in I\}$ is a simplex in
$\R^R$, without an additional origin.

\begin{theorem}[Postnikov]\label{thm:Postnikov}
For nonnegative integers $y_0,\ldots,y_m$,
\begin{equation}\label{eq:Postnikov}
 \#\left(\left(\sum_{i=0}^m y_i\Delta_{I_i}\right)\cap\Z^R\right)
 =\sum_{\alpha\text{ draconian}}
     \binom{y_0+\alpha_0}{\alpha_0}
     \prod_{i=1}^m\binom{y_i+\alpha_i-1}{\alpha_i}.
\end{equation}
\end{theorem}
This is the untrimmed form of \cite[Theorem~11.3]{Postnikov09}, obtained by
increasing by one the coefficient of the full simplex. It is also the form
used in \cite[Theorem~2.1]{AC26}. Repeated neighborhoods are permitted, as are
zero coefficients. The offset in the $i=0$ factor is essential.

All binomial coefficients in polynomial identities mean
\begin{equation}\label{eq:polynomial-binom}
 \binom zk=\frac{z(z-1)\cdots(z-k+1)}{k!}\quad(k\in\N),\qquad
 \binom z0=1.
\end{equation}
In particular, setting a coefficient $y_i=0$ annihilates its factors with
$\alpha_i>0$ when $i\ne0$, but setting $y_0=0$ does not annihilate the
$\alpha_0$ factor.

\section{Weighted preorder polytopes}\label{sec:weighted}

\subsection{A Minkowski-sum description with independent capacities}
For $S\subseteq E$, write
$\ds_S=\conv(\{\zero\}\cup\{\mathbf e_e:e\in S\})$ in $\R^E$;
thus $\ds_\varnothing=\{\zero\}$. The bold simplex includes the origin,
unlike $\Delta_I$ in \Cref{thm:Postnikov}.

For $\bvec\in\R_{\geq0}^E$ and $s\geq0$, define
\begin{equation}\label{eq:weighted-def}
 \QQ_\tau(\bvec;s)=
 \{\xvec\geq\zero:x(I)\leq b(I)+s
                    \text{ for all }\varnothing\ne I\in\JJ(\tau)\}.
\end{equation}
The uniform family is $\QQ_\tau(r,s)=\QQ_\tau(r\one;s)$.

\begin{proposition}\label{prop:minkowski}
For arbitrary nonnegative real parameters,
\begin{equation}\label{eq:weighted-minkowski}
 \QQ_\tau(\bvec;s)=\sum_{e\in E}b_e\ds_{U_\tau(e)}+s\ds_E.
\end{equation}
Consequently, this is a lattice polytope when the parameters are integers.
\end{proposition}
\begin{proof}
Denote the right-hand side by $M$. If $I$ is an ideal, the summand indexed by
$e$ can contribute to a coordinate in $I$ only if $e\in I$. Its total
contribution to $x(I)$ is at most $b_e$. The last summand contributes at most
$s$. Therefore $M\subseteq\QQ_\tau(\bvec;s)$.

For the reverse containment we compare linear functionals. Both sets are
closed under decreasing nonnegative coordinates. Hence it suffices to
consider a functional with nonnegative coefficients $\cvec$; for a general
functional, replace its coefficients by their positive parts and note that
the support function of $M$ does not change.

For $t>0$ let $S_t=\{f\in E:c_f\geq t\}$. For a feasible point,
whenever $S_t\ne\varnothing$,
\[
 x(S_t)\leq x(D_\tau(S_t))\leq b(D_\tau(S_t))+s.
\]
Integrate over $0<t\leq\max_f c_f$. The elementary layer-cake identity gives
\begin{align*}
 \cvec\cdot\xvec
 &=\int_0^{\max c_f}x(S_t)\,dt\\
 &\leq\sum_{e\in E}b_e\max_{f\in U_\tau(e)}c_f+s\max_{f\in E}c_f.
\end{align*}
The last expression is exactly the maximum of $\cvec\cdot\xvec$ on $M$,
since each simplex summand can choose a coordinate attaining its own maximum.
Thus no linear functional separates a feasible point from $M$; compact
convex separation proves the desired containment. Boundedness follows from
the ideal $E$.

For integer parameters the right-hand side is a Minkowski sum of lattice
polytopes, hence is the convex hull of integer sums of their vertices.
It is therefore a lattice polytope.
\end{proof}

\begin{remark}\label{rem:ideals}
At $s=0$ some inequalities for disconnected ideals are sums of inequalities
for their components. This is not a license to omit those ideals when
$s>0$: summing component bounds would add the slack $s$ several times.
We keep every nonempty ideal throughout.
\end{remark}

\subsection{The multivariate lattice enumerator}
Introduce the graph with left and right copies of $E\sqcup\{0\}$ and
neighborhoods
\[
 I_0=E\sqcup\{0\},\qquad I_e=U_\tau(e)\sqcup\{0\}\quad(e\in E).
\]
The projection dropping coordinate~$0$ identifies the corresponding
Minkowski sum, with weights $y_0=s$ and $y_e=b_e$, with
\eqref{eq:weighted-minkowski}. For integer parameters, its inverse on the
fixed-sum hyperplane is
$x_0=\sumv{\bvec}+s-\sumv{\xvec}$, so the projection preserves the affine
integer lattices.

The following identification is \cite[Lemma~4.1]{AC26}; we recall its
proof to make the dual orientation explicit.

\begin{lemma}\label{lem:draconian-preorder}
A vector $(a_0,(a_e)_{e\in E})$ is draconian for this graph if and only if
$\avec\in P_{\tau^*}$ and $a_0=n-\sumv{\avec}$.
\end{lemma}
\begin{proof}
The total-sum condition gives the asserted value of $a_0$. A subset of left
vertices containing $0$ has all $n+1$ right vertices as neighbors, so its
inequality follows from nonnegativity and the total sum $n$. For a nonempty
subset $A\subseteq E$, the neighborhood inequality is
\[
 a(A)\leq\card{U_\tau(A)}.
\]
If $A$ itself is a filter, this is $a(A)\leq\card A$. Conversely, the
inequalities for all filters imply the inequality for an arbitrary $A$,
since $a(A)\leq a(U_\tau(A))\leq\card{U_\tau(A)}$. Filters of $\tau$
are ideals of $\tau^*$. The condition for the filter $E$ also guarantees
$a_0\geq0$.
\end{proof}

\begin{proposition}\label{prop:F}
There is a polynomial $\FF_\tau(\bvec,s)$ of total degree $n$ such that
\[
 \FF_\tau(\bvec,s)=\#(\QQ_\tau(\bvec;s)\cap\Z^E)
 \qquad(\bvec\in\N^E,\ s\in\N).
\]
It is given by
\begin{equation}\label{eq:rising}
 \boxed{\FF_\tau(\bvec,s)=
 \sum_{\avec\in P_{\tau^*}}
   \binom{s+n-\sumv{\avec}}{n-\sumv{\avec}}
   \prod_{e\in E}\binom{b_e+a_e-1}{a_e}.}
\end{equation}
In particular, $\EE_\tau(u,v)=\FF_\tau(u\one,v)$.
\end{proposition}
\begin{proof}
Apply \Cref{thm:Postnikov} and \Cref{lem:draconian-preorder}. Every summand
has total degree at most $n$. The degree is exactly $n$, because setting
$\bvec=\zero$ leaves
$\FF_\tau(\zero,s)=\binom{s+n}{n}$, the lattice count of $s\ds_E$.
Uniqueness follows because polynomials agreeing on $\N^{n+1}$ agree
identically; this can be proved one variable at a time.
\end{proof}

\section{Reciprocity by interior translation}\label{sec:reciprocity}

\begin{lemma}[Interior translation]\label{lem:translation}
Suppose $b_e\geq1$ for every $e\in E$ and $s\geq1$, all integers. Then
$\QQ_\tau(\bvec;s)$ is full dimensional, and translation by $-\one$ gives
a bijection
\begin{equation}\label{eq:weighted-translation}
 \Int\QQ_\tau(\bvec;s)\cap\Z^E
 \ \longleftrightarrow\
 \QQ_\tau(\bvec-\one;s-1)\cap\Z^E.
\end{equation}
\end{lemma}
\begin{proof}
The point $\one$ satisfies all coordinate inequalities strictly and all
nonempty ideal inequalities strictly, since
$\card I<b(I)+s$. There are finitely many inequalities, so a sufficiently
small Euclidean ball about $\one$ lies in the polytope. This proves full
dimension.

For a full-dimensional polytope described by finitely many inequalities with
nonzero normals, a point is interior exactly when every defining inequality
is strict. Redundant inequalities cause no problem: equality in any valid
inequality with nonzero normal prevents a full ball around the point. Thus an
integer point $\xvec$ is interior if and only if
\[
 x_e\geq1\quad(e\in E),\qquad
 x(I)\leq b(I)+s-1\quad(\varnothing\ne I\in\JJ(\tau)).
\]
Put $\yvec=\xvec-\one$. The coordinate conditions become $y_e\geq0$, and
\[
 y(I)\leq b(I)-\card I+s-1=(\bvec-\one)(I)+(s-1).
\]
These are exactly the inequalities on the right of
\eqref{eq:weighted-translation}. All transformations are reversible over
$\Z^E$.
\end{proof}

\begin{theorem}[Multivariate reciprocity]\label{thm:reciprocity}
For every finite preorder on $n\geq1$ elements,
\begin{equation}\label{eq:weighted-reciprocity}
 \boxed{\FF_\tau(-\bvec,-s)=(-1)^n\FF_\tau(\bvec-\one,s-1)}
\end{equation}
as an identity in $\Q[(b_e)_{e\in E},s]$.
\end{theorem}
\begin{proof}
First fix positive integer parameters $\bvec,s$. The polytope
$K=\QQ_\tau(\bvec;s)$ is a full-dimensional lattice polytope by
\Cref{prop:minkowski,lem:translation}. For every nonnegative integer $t$,
\[
 tK=\QQ_\tau(t\bvec;ts),\qquad
 L_K(t)=\FF_\tau(t\bvec,ts).
\]
Both sides of the second equality are polynomials in $t$, so they agree
also at $t=-1$. Ehrhart--Macdonald reciprocity therefore gives
\begin{align*}
 \FF_\tau(-\bvec,-s)
 &=L_K(-1)\\
 &=(-1)^n\#(\Int K\cap\Z^E)\\
 &=(-1)^n\FF_\tau(\bvec-\one,s-1),
\end{align*}
where the last equality is \Cref{lem:translation}.

The difference between the two sides of
\eqref{eq:weighted-reciprocity} is a polynomial vanishing on every tuple of
positive integers. Fix all variables but one: the remaining univariate
polynomial has infinitely many roots. Repeat this observation for the
coefficient polynomials in the remaining variables. The difference is zero
identically.
\end{proof}

\begin{corollary}[Conjecture 7.2 of \cite{AC26}]\label{cor:double}
For all indeterminates $u,v$,
\[
 \EE_\tau(-u,-v)=(-1)^n\EE_\tau(u-1,v-1).
\]
\end{corollary}
\begin{proof}
Set $b_e=u$ for every $e$ in \Cref{thm:reciprocity}.
\end{proof}

\subsection{Why the polynomial step matters}
We intentionally proved the lattice bijection only when every parameter is
positive. For example, $\QQ_\tau(\zero;0)$ is a point, not a full-dimensional
polytope. The identity at zero or negative parameters follows from the
polynomial continuation, not from applying an interior formula to a
lower-dimensional or nonexistent feasible set. Similarly, no assertion
about the interior of a Minkowski sum being a sum of interiors is needed.
The direct inequalities avoid that issue entirely.

\begin{corollary}[Centered parity]\label{cor:parity}
Define
\[
 H_\tau(\bvec,s)=\FF_\tau(\bvec-\tfrac12\one,s-\tfrac12).
\]
Then $H_\tau(-\bvec,-s)=(-1)^n H_\tau(\bvec,s)$. Consequently, only
monomials whose total degree has the same parity as $n$ occur in $H_\tau$.
For the double polynomial, if $E_j(u,v)$ is its homogeneous part of degree
$j$, then
\begin{equation}\label{eq:coeff-identity}
 E_{n-1}(u,v)=\tfrac12(\partial_u+\partial_v)E_n(u,v).
\end{equation}
\end{corollary}
\begin{proof}
Substitute $\bvec+\tfrac12\one,s+\tfrac12$ into
\eqref{eq:weighted-reciprocity}. A monomial of degree $j$ changes by
$(-1)^j$ under simultaneous negation. For the last identity, write
$\EE_\tau(u,v)=H(u+\tfrac12,v+\tfrac12)$, where $H$ has no homogeneous
part of degree $n-1$. The degree-$n-1$ terms come solely from translating
its degree-$n$ part, giving \eqref{eq:coeff-identity}.
\end{proof}

\section{A second binomial formula and the word conjecture}\label{sec:words}

Athanasiadis and Chapoton already explain in \cite[Remark~7.3]{AC26}
that Conjecture~7.2 implies their word conjecture, and in
\cite[Remark~4.10(a)]{AC26} that the word conjecture implies the ordinary
zeta evaluation. We include these deductions in full, with an independent
proof of the elementary word identity. The multivariate falling-binomial
formula below strengthens the uniform specialization used there.

\subsection{The falling-binomial form}

\begin{theorem}\label{thm:falling}
For every finite preorder,
\begin{equation}\label{eq:falling}
 \boxed{\FF_\tau(\bvec,s)=
 \sum_{\avec\in P_{\tau^*}}
       \binom{s}{n-\sumv{\avec}}
       \prod_{e\in E}\binom{b_e+1}{a_e}.}
\end{equation}
In particular,
\begin{equation}\label{eq:bases-formula}
 \FF_\tau(\bvec,0)=\sum_{\avec\in B_{\tau^*}}
                      \prod_{e\in E}\binom{b_e+1}{a_e},
 \qquad
 L_{\QQ_\tau}(m)=\sum_{\avec\in B_{\tau^*}}
                      \prod_{e\in E}\binom{m+1}{a_e}.
\end{equation}
\end{theorem}
\begin{proof}
By \Cref{thm:reciprocity},
$\FF_\tau(\bvec,s)=(-1)^n\FF_\tau(-\bvec-\one,-s-1)$.
Apply \eqref{eq:rising} to the latter expression. Put
$a_0=n-\sumv{\avec}$. The polynomial identities
\[
 \binom{-s-1+a_0}{a_0}=(-1)^{a_0}\binom{s}{a_0},\qquad
 \binom{-b_e+a_e-2}{a_e}=(-1)^{a_e}\binom{b_e+1}{a_e}
\]
show that the sign of each summand is
$(-1)^{a_0+\sum_e a_e}=(-1)^n$. It cancels the external sign. This yields
\eqref{eq:falling}. At $s=0$, the factor $\binom0{a_0}$ is nonzero exactly
when $a_0=0$. The surviving vectors form $B_{\tau^*}$ by
\Cref{lem:maximal}.
\end{proof}

For nonnegative parameters, \eqref{eq:falling} has a simple counting
interpretation. Take disjoint sets of sizes $s$ and $b_e+1$ for $e\in E$.
Choose a total of $n$ objects, with the vector of choices from the
$E$-indexed sets lying in $P_{\tau^*}$. Their number is
$\FF_\tau(\bvec,s)$. This interpretation counts sets without repetitions,
whereas \eqref{eq:rising} uses selections with repetitions. The equality is
an enumerative consequence of reciprocity; we do not assert a canonical
bijection between these two models.

\subsection{A self-contained multiset-word identity}
Fix any total order on $E$, independently of the preorder. If
$\sumv{\avec}=n$, let $\mathfrak S_\avec$ be the words of content $\avec$.
For $w=w_1\cdots w_n$ define
$\des(w)=\#\{i\in\{1,\ldots,n-1\}:w_i>w_{i+1}\}$.

\begin{lemma}\label{lem:word}
For any $\avec\in\N^E$ of coordinate sum $n\geq1$,
\begin{equation}\label{eq:word-identity}
 \sum_{m\geq0}\left(\prod_{e\in E}\binom{m+1}{a_e}\right)t^m
 =\frac{\displaystyle\sum_{w\in\mathfrak S_\avec}
                          t^{n-1-\des(w)}}{(1-t)^{n+1}}.
\end{equation}
\end{lemma}
\begin{proof}
For each letter $e$, choose $a_e$ distinct levels from
$\{0,1,\ldots,m\}$. At each level, list its selected letters in strictly
\emph{decreasing} order, then concatenate the levels in increasing order.
This produces a word $w$ of content $\avec$, together with a weakly
increasing sequence of levels attached to its positions.

Two adjacent positions may share a level only when they are a strict
descent in $w$. Thus a strict increase of levels is required at precisely
the $d=n-1-\des(w)$ non-descent positions. Conversely, such a level
assignment reconstructs the subsets of levels: within an equal-level block
the letters are strictly decreasing, so no letter occurs there twice.

Subtract, from each level, the number of required strict increases before
that position. This leaves a weakly increasing $n$-tuple with entries in
$\{0,\ldots,m-d\}$. Its number is
$\binom{m-d+n}{n}$ when $m\geq d$ and zero otherwise. Summing over $m$
yields
\[
 \sum_{m\geq d}\binom{m-d+n}{n}t^m
 =\frac{t^d}{(1-t)^{n+1}}.
\]
Finally sum over $w\in\mathfrak S_\avec$.
\end{proof}

\begin{theorem}[Conjecture 4.4 of \cite{AC26}]\label{thm:hstar}
Let $\WW_\tau$ be the words of length $n$ on the totally ordered alphabet
$E$ for which every order ideal $I$ of $\tau$ contributes at least
$\card I$ letters. Then
\begin{equation}\label{eq:hstar-word}
 \boxed{h^*(\QQ_\tau,t)=\sum_{w\in\WW_\tau}t^{n-1-\des(w)}.}
\end{equation}
\end{theorem}
\begin{proof}
A content vector $\avec$ of total sum $n$ belongs to $B_{\tau^*}$ if and
only if $a(F)\leq\card F$ for every filter $F$ of $\tau$. Complementing
$F$ converts these inequalities into $a(I)\geq\card I$ for every ideal
$I$. Hence
$\WW_\tau=\bigsqcup_{\avec\in B_{\tau^*}}\mathfrak S_\avec$.
Sum \eqref{eq:word-identity} over $B_{\tau^*}$, use
\eqref{eq:bases-formula}, and compare the result with
\eqref{eq:hstar-def}.
\end{proof}

\subsection{Extremal coefficients, volume, and zeta evaluation}
The proof immediately gives
\begin{equation}\label{eq:extremal}
 h_0^*(\QQ_\tau)=1,\qquad h_n^*(\QQ_\tau)=0,\qquad
 h_{n-1}^*(\QQ_\tau)=\#B_{\tau^*}.
\end{equation}
For the constant coefficient, a word with $n-1$ descents is the unique
strictly decreasing listing of all $n$ letters. For the coefficient at
$n-1$, each admissible content has exactly one weakly increasing word.
The total number of admissible words is
\begin{equation}\label{eq:volume}
 n!\Vol(\QQ_\tau)=\sum_{\avec\in B_{\tau^*}}
                          \frac{n!}{\prod_{e\in E}a_e!}.
\end{equation}
This follows either from the leading coefficient in
\eqref{eq:bases-formula} or by setting $t=1$ in the numerator.

The absence of interior lattice points in $\QQ_\tau$ can also be seen
directly: such a point would have every coordinate at least one, while
strictness of the inequality for $E$ would force their sum below $n$.
The point $\one$ is interior to $2\QQ_\tau$. Thus the smallest positive
dilation with an interior lattice point is exactly two, consistently with
$\deg h^*=n-1$.

For $k\geq1$, count the $k$-multichains in $P_\tau$ by their top element.
Each coordinate gives a weakly increasing sequence ending at $a_e$, so
\begin{equation}\label{eq:zeta}
 Z(P_\tau,k+1)=\sum_{\avec\in P_\tau}
                      \prod_{e\in E}\binom{k+a_e-1}{a_e}.
\end{equation}
This equation defines its polynomial continuation. Comparing with
\eqref{eq:rising} at $s=0$ gives
$Z(P_\tau,t)=L_{\QQ_{\tau^*}}(t-1)$.

\begin{corollary}[Ordinary part of Conjecture 4.9]\label{cor:zeta}
For every preorder on $n$ elements,
\begin{equation}\label{eq:zeta-minus-one}
 Z(P_\tau,-1)=(-1)^n\#\Max(P_\tau).
\end{equation}
Equivalently,
$\#(\Int(2\QQ_{\tau^*})\cap\Z^E)=\#\Max(P_\tau)$.
\end{corollary}
\begin{proof}
Apply \eqref{eq:bases-formula} to $\tau^*$ at $m=-2$. Since
$\binom{-1}{a_e}=(-1)^{a_e}$ and every base has sum $n$,
\[
 Z(P_\tau,-1)=L_{\QQ_{\tau^*}}(-2)
 =\sum_{\avec\in B_\tau}\prod_e\binom{-1}{a_e}
 =(-1)^n\#B_\tau.
\]
Use \Cref{lem:maximal} for the first assertion and
\eqref{eq:EM} for the second.
\end{proof}
The source's refined conjecture assigns powers of $q$ to covers of maximal
elements. None of the preceding equalities transports that statistic. The
ordinary specialization therefore must not be reported as a proof of the
full refined conjecture.

\section{Colored arrays and matrix Ehrhart--zeta duality}\label{sec:matrix}

The preceding reciprocity proof is independent of this section. Here the
key idea is to replicate \emph{both} color classes in the preorder graph.
The two replication counts need not be equal, or even constant over $E$.

\subsection{The heterogeneous enumeration problem}
For $\pvec,\qvec\in\N^E$, let
\[
 R_{\qvec}=\{(e,j):e\in E,\ 1\leq j\leq q_e\}.
\]
Define $C_\tau(\pvec,\qvec)$ as the number of arrays
$(x_{e,j})_{(e,j)\in R_{\qvec}}\in\N^{R_{\qvec}}$ satisfying
\begin{equation}\label{eq:colored-constraints}
 \sum_{e\in I}\sum_{j=1}^{q_e}x_{e,j}\leq p(I)
 \qquad(I\in\JJ(\tau)).
\end{equation}
Empty groups of coordinates are permitted. The empty array contributes one
when all $q_e=0$.

Aggregating by $a_e=\sum_jx_{e,j}$ gives the explicit finite sum
\begin{equation}\label{eq:colored-sum}
 C_\tau(\pvec,\qvec)=
 \sum_{\substack{\avec\in\N^E\\a(I)\leq p(I)\ (I\in\JJ(\tau))}}
             \prod_{e\in E}\binom{q_e+a_e-1}{a_e}.
\end{equation}
For $q_e=0$, the factor is one for $a_e=0$ and zero for $a_e>0$, as it
should be. The ideal $E$ bounds $\sumv{\avec}$ by $\sumv{\pvec}$, so the
sum is finite.

\begin{theorem}[Heterogeneous array duality]\label{thm:colored}
For every finite preorder and all nonnegative integer vectors
$\pvec,\qvec$,
\begin{equation}\label{eq:colored-duality}
 \boxed{C_\tau(\pvec,\qvec)=C_{\tau^*}(\qvec,\pvec).}
\end{equation}
\end{theorem}
We prove this by identifying the two sides with the lattice count and the
draconian count of one bipartite graph. The graph-duality phenomenon itself
is classical in Postnikov's theory \cite[Corollary~11.8]{Postnikov09}; the
point here is the exact encoding of the preorder constraints and of two
independently varying multiplicity vectors.

\subsection{The copied graph and its polytope}
Let the left vertices be
\[
 L=\{0_L\}\sqcup\{(e,i)_L:e\in E,\ 1\leq i\leq p_e\}
\]
and the right vertices be
\[
 R=\{0_R\}\sqcup\{(f,j)_R:f\in E,\ 1\leq j\leq q_f\}.
\]
Join $0_L$ to every right vertex and every left vertex to $0_R$.
For the other vertices, put an edge precisely when
\begin{equation}\label{eq:copy-edge}
 (e,i)_L\sim(f,j)_R\quad\Longleftrightarrow\quad e\leq_\tau f.
\end{equation}
Denote this graph by $G_\tau(\pvec,\qvec)$. Its left neighborhoods are
nonempty because they all contain $0_R$, and the neighborhood of $0_L$
is all of $R$.

Let $K$ be the projection, dropping coordinate $0_R$, of
\begin{equation}\label{eq:copy-minkowski}
 \sum_{\lambda\in L\setminus\{0_L\}}\Delta_{N(\lambda)}.
\end{equation}
Equivalently, with
$S_e=\{(f,j):e\leq_\tau f,\ 1\leq j\leq q_f\}$,
\[
 K=\sum_{e\in E}p_e\ds_{S_e}\subseteq\R^{R_{\qvec}}.
\]
The coefficient of the full simplex indexed by $0_L$ is zero.

\begin{lemma}\label{lem:colored-polytope}
The polytope $K$ is precisely the nonnegative real solution set of
\eqref{eq:colored-constraints}. Hence $\#(K\cap\Z^{R_{\qvec}})
=C_\tau(\pvec,\qvec)$.
\end{lemma}
\begin{proof}
Any summand associated with $e$ can contribute to coordinates with labels
in an ideal $I$ only if $e\in I$. This proves that points in $K$ satisfy
\eqref{eq:colored-constraints}.

For the converse, repeat the support-function proof of
\Cref{prop:minkowski}, now on copied coordinates. Given nonnegative
coefficients $c_{f,j}$, take a level set of copied coordinates and let $T_t$
be the set of their original labels. Its downward closure $D_\tau(T_t)$ is
an ideal. The sum of the selected coordinates is at most the sum of all
coordinates with labels in this ideal, and hence at most $p(D_\tau(T_t))$.
Integration over levels bounds the functional by
\[
 \sum_{e\in E}p_e\max\bigl(\{c_{f,j}:e\leq_\tau f,\ 1\leq j\leq q_f\}
                             \cup\{0\}\bigr).
\]
This is the support function of $K$. The extra $0$ handles empty
neighborhoods after projection. Downward closure handles negative
coefficients, and convex separation completes the proof, just as before.
\end{proof}

The unprojected polytope in \eqref{eq:copy-minkowski} lies in the hyperplane
with coordinate sum $\sumv{\pvec}$. Projection preserves integer points:
the missing coordinate is
$\sumv{\pvec}-\sum_{(f,j)\in R_{\qvec}}x_{f,j}$.
Thus \Cref{thm:Postnikov} can count $K$.

\subsection{Reading the draconian inequalities on the other side}
Apply \eqref{eq:Postnikov} with coefficient zero at $0_L$ and coefficient
one at every other left vertex. Every summand in that formula equals one:
\[
 \binom{\alpha_0}{\alpha_0}=1,\qquad
 \binom{1+\alpha_\lambda-1}{\alpha_\lambda}=1.
\]
Consequently,
\begin{equation}\label{eq:count-draconian}
 C_\tau(\pvec,\qvec)=\#\{G_\tau(\pvec,\qvec)\text{-draconian vectors}\}.
\end{equation}

Let $Q=\sumv{\qvec}$, so that the right color class has $Q+1$ vertices.
A draconian vector has coordinates
$\alpha_0$ and $\alpha_{e,i}$, with
\begin{equation}\label{eq:copy-total}
 \alpha_0+\sum_{e\in E}\sum_{i=1}^{p_e}\alpha_{e,i}=Q.
\end{equation}
For a nonempty set $A$ of non-root left vertices, let $S$ be its set of
original labels. Its neighborhood has size
$1+q(U_\tau(S))$, so its inequality is
\begin{equation}\label{eq:copy-subset}
 \alpha(A)\leq q(U_\tau(S)).
\end{equation}
These inequalities are equivalent to
\begin{equation}\label{eq:copy-filter}
 \sum_{e\in F}\sum_{i=1}^{p_e}\alpha_{e,i}\leq q(F)
 \qquad(F\text{ an order filter of }\tau).
\end{equation}
Here is a detailed check, including zero multiplicities. Given a filter $F$,
take $A$ to be all the existing left copies with labels in $F$. If $A$ is
empty, \eqref{eq:copy-filter} is automatic. Otherwise its label set $S$
satisfies $U_\tau(S)\subseteq F$, and
\eqref{eq:copy-subset} implies \eqref{eq:copy-filter}. Conversely, for an
arbitrary nonempty $A$, its coordinates are a subset of all left copies
whose labels lie in the filter $U_\tau(S)$. Nonnegativity and
\eqref{eq:copy-filter} imply \eqref{eq:copy-subset}.

Sets of left vertices containing $0_L$ have the whole right color class as
neighborhood, so their inequalities follow from \eqref{eq:copy-total}.
The filter $E$ in \eqref{eq:copy-filter} makes
\[
 \alpha_0=Q-\sum_{e,i}\alpha_{e,i}
\]
nonnegative and determines it uniquely. Therefore the draconian vectors
are in bijection with nonnegative arrays on the $p_e$ left copies satisfying
\eqref{eq:copy-filter}. Since filters of $\tau$ are ideals of $\tau^*$,
these arrays are exactly those counted by
$C_{\tau^*}(\qvec,\pvec)$. Combined with
\eqref{eq:count-draconian}, this proves \Cref{thm:colored}.

\subsection{Returning to multichains}
Let $\nabla_\tau(k,\ell)$ denote the number of $k$-multichains in the
coordinatewise ordered lattice points of $\ell\QQ_\tau$. A multichain
has $k$ entries; the empty multichain is counted once when $k=0$.

For $k\geq1$, convert a multichain
$\xvec^{(1)}\leq\cdots\leq\xvec^{(k)}$ into its coordinate increments:
\[
 y_{e,1}=x_e^{(1)},\qquad
 y_{e,j}=x_e^{(j)}-x_e^{(j-1)}\quad(2\leq j\leq k).
\]
These are nonnegative. Since $\ell\QQ_\tau$ is a downset, all the
multichain entries are feasible exactly when its last entry is feasible.
The latter condition says
\[
 \sum_{e\in I}\sum_{j=1}^k y_{e,j}\leq\ell\card I.
\]
Partial sums provide the inverse map. Thus, including $k=0$,
\begin{equation}\label{eq:matrix-C}
 \nabla_\tau(k,\ell)=C_\tau(\ell\one,k\one).
\end{equation}

\begin{corollary}[Conjecture 4.2 of \cite{AC26}]\label{cor:matrix}
For every preorder and all $k,\ell\in\N$,
\begin{equation}\label{eq:matrix-duality}
 \boxed{\nabla_\tau(k,\ell)=\nabla_{\tau^*}(\ell,k).}
\end{equation}
\end{corollary}
\begin{proof}
Use \eqref{eq:matrix-C}, then \Cref{thm:colored}, then
\eqref{eq:matrix-C} again with $\tau^*$.
\end{proof}

When either parameter is zero, both sides are one. No argument with
fractional dilations of a copied preorder is involved. All coefficients
in the copied Minkowski sum are zero or one, even though its ambient
dimension varies with $\pvec,\qvec$.

Self-duality of the preorder is sufficient for symmetry of $\nabla_\tau$,
but symmetry is not asserted in general. Nor does
\eqref{eq:colored-duality} say that $C_\tau$ is jointly a polynomial of
bounded degree in the entries of both vectors. Already for a singleton,
$C(p,q)=\binom{p+q}{p}$.

\section{Worked examples}\label{sec:examples}

\subsection{Antichains and a single equivalence class}
For an $n$-element antichain, $\QQ_\tau=[0,1]^n$ and
$P_\tau=\{0,1\}^n$. Formula \eqref{eq:falling} yields
\begin{equation}\label{eq:antichain-double}
 \EE_\tau(r,s)=\sum_{j=0}^n\binom nj\binom{s}{n-j}(r+1)^j.
\end{equation}
Indeed, at a fixed support of size $j$, the only vector in $P_\tau$ with
that support has all its nonzero entries equal to one. The heterogeneous
array count factors coordinate by coordinate:
\[
 C_\tau(\pvec,\qvec)=\prod_{e\in E}\binom{p_e+q_e}{p_e},\qquad
 \nabla_\tau(k,\ell)=\binom{k+\ell}{k}^{\!n}.
\]
Both dualities are visibly symmetric here.

At the other extreme, suppose every two elements of $E$ are equivalent.
There is only one nonempty ideal, namely $E$. Then
\[
 \QQ_\tau(\bvec;s)=(\sumv{\bvec}+s)\ds_E,\qquad
 \FF_\tau(\bvec,s)=\binom{\sumv{\bvec}+s+n}{n},
\]
and
\[
 C_\tau(\pvec,\qvec)=\binom{\sumv{\pvec}+\sumv{\qvec}}{\sumv{\pvec}},
 \qquad
 \nabla_\tau(k,\ell)=\binom{n(k+\ell)}{nk}.
\]
These examples confirm that the cardinality $n$ counts elements, not
preorder equivalence classes.

\subsection{The two-element chain}
Let $1<_{\tau}2$. Its inequalities are
$x_1\leq r+s$ and $x_1+x_2\leq2r+s$, with nonnegative coordinates.
Direct summation gives
\begin{align*}
 \EE_\tau(r,s)
 &=\sum_{x_1=0}^{r+s}(2r+s-x_1+1)\\
 &=\frac{(r+s+1)(3r+s+2)}2.
\end{align*}
After centering,
\[
 \EE_\tau(u-\tfrac12,v-\tfrac12)=\frac{(u+v)(3u+v)}2,
\]
which has the even parity predicted by \Cref{cor:parity}. At $s=0$,
$L_{\QQ_\tau}(r)=(r+1)(3r+2)/2$ and $h^*(\QQ_\tau,t)=1+2t$.

\subsection{A non-self-dual pair of three-element posets}
Let $V$ be the poset with $1<2$ and $1<3$, and let $\Lambda=V^*$.
Its two Ehrhart polynomials are
\begin{align}
 L_{\QQ_V}(m)&=\frac{(m+1)(13m^2+17m+6)}6,\label{eq:forkL}\\
 L_{\QQ_\Lambda}(m)&=(m+1)^2(2m+1).
\end{align}
Consequently,
\[
 h^*(\QQ_V,t)=1+8t+4t^2,\qquad
 h^*(\QQ_\Lambda,t)=1+8t+3t^2.
\]
Both posets of lattice points have $12$ elements, but their Ehrhart and
$h^*$-polynomials differ. This is not a conflict with matrix duality:
that theorem exchanges the two parameters as well as the preorder.

In fact the initial matrices, with rows indexed by $k=0,1,2,3$ and columns
by $\ell=0,1,2,3$, are
\begin{equation}\label{eq:fork-matrices}
 \nabla_V=
 \begin{pmatrix}
 1&1&1&1\\
 1&12&46&116\\
 1&45&436&2235\\
 1&112&2200&20010
 \end{pmatrix},\qquad
 \nabla_\Lambda=
 \begin{pmatrix}
 1&1&1&1\\
 1&12&45&112\\
 1&46&436&2200\\
 1&116&2235&20010
 \end{pmatrix}.
\end{equation}
They are transposes, and neither is symmetric. A heterogeneous example is
\[
 C_V((1,2,1),(2,1,3))
 =C_\Lambda((2,1,3),(1,2,1))=58.
\]

For the double polynomial of $V$, direct enumeration or
\eqref{eq:rising} gives
\[
 \EE_V(u,v)=\frac{(u+v+1)(13u^2+8uv+17u+v^2+5v+6)}6.
\]
Its centered form is particularly transparent:
\[
 \EE_V(u-\tfrac12,v-\tfrac12)
 =\frac{(u+v)(26u^2+16uv+2v^2+1)}{12}.
\]
Only degrees three and one occur.

There are three maximal elements of $P_V$:
$(0,1,2),(0,2,1),(1,1,1)$. There are four in $P_\Lambda$:
$(1,1,1),(2,0,1),(2,1,0),(3,0,0)$. Thus the respective leading nonzero
$h^*$ coefficients are four and three, as
\eqref{eq:extremal} predicts. The corresponding zeta evaluations are
$Z(P_V,-1)=-3$ and $Z(P_\Lambda,-1)=-4$.

\subsection{A mixed preorder with an equivalence class}
On $E=\{a,b,c,d,e\}$ let $a$ and $c$ be equivalent, let both lie below
$e$, and impose $b<d$ and $b<e$. No further comparisons are imposed except
those forced by reflexivity and transitivity. The exact enumeration gives
\[
 \#P_\tau=92,\qquad \#\Max(P_\tau)=18,
\]
\[
 L_{\QQ_\tau}(m)
 =\frac{(m+1)^2(2m+1)(19m^2+21m+6)}6,
\]
\[
 h^*(\QQ_\tau,t)=1+86t+393t^2+260t^3+20t^4.
\]
The maximal-element count of the dual is $20$, while
$Z(P_\tau,-1)=-18$. This provides a check on both the dual orientation and
the distinction between the two maximal-element counts. The entire double
polynomial, in exact Newton-basis form, is supplied in the accompanying
JSON data.

\section{Exact verification and reproducible artifacts}\label{sec:verification}

\subsection{Independent enumeration routes}
The program \texttt{code/verify.py} uses only the Python standard library
and exact integers. It exhausts all reflexive binary relations on
$\{0,\ldots,n-1\}$ and keeps precisely those satisfying transitivity.
Thus it includes all labeled preorders, not merely naturally labeled
partial orders or representatives of a few Hasse-diagram shapes.

Lattice points are generated directly from nonnegative vectors whose total
sum is bounded by the ideal $E$, then tested against every nonempty ideal.
The program compares this count with both binomial formulas. The word
formula is checked by enumerating all $n^n$ words and testing the lower-ideal
conditions directly. Maximal elements are found by pairwise coordinatewise
comparison, independently of the coordinate-sum criterion.

For matrix and heterogeneous duality, the program evaluates
\eqref{eq:colored-sum} separately for $\tau$ and $\tau^*$. This is not an
implementation of the graph argument. For interior translation it compares
actual finite sets, not only their cardinalities.

\subsection{Polynomial checking rather than isolated evaluations}
For each preorder, the double polynomial is reconstructed from direct
lattice counts at
\[
 (r,s)\in\N^2,\qquad r+s\leq n.
\]
In the Newton basis it is
\begin{equation}\label{eq:newton}
 \EE_\tau(u,v)=\sum_{i+j\leq n}c_{i,j}\binom ui\binom vj,
\end{equation}
where
\[
 c_{i,j}=\sum_{a=0}^i\sum_{b=0}^j
       (-1)^{i+j-a-b}\binom ia\binom jb\EE_\tau(a,b).
\]
This triangular grid is unisolvent for polynomials of total degree at most
$n$: order its points by total degree and note that a Newton basis element
vanishes at earlier grid points and equals one at its own point. The
reciprocity difference, also of degree at most $n$, is tested on the
translated triangular grid $(r+1,s+1)$. For each fixed preorder, these
checks certify that difference as a polynomial, given the proved degree
bound. They do not replace the proof for arbitrary preorders.

\subsection{Recorded run}
The supplied run covers $1,4,29,355$ labeled preorders of sizes
$1,2,3,4$, respectively: $389$ in total. Its results are:
\begin{center}
\begin{tabularx}{\textwidth}{@{}Xr@{}}
\toprule
Check & Number \\
\midrule
Direct lattice-count grid points, compared with both formulas & 5,642 \\
Unisolvent-grid reciprocity tests & 5,642 \\
Interior translation: equality of actual point sets & 1,556 \\
Maximal-element characterizations & 389 \\
Full $h^*$ vectors versus direct word histograms & 389 \\
Ordinary zeta evaluations at $-1$ & 389 \\
Matrix duality, $0\leq k,\ell\leq3$ & 6,224 \\
Heterogeneous array duality cases & 3,112 \\
Heterogeneous reciprocity cases & 3,112 \\
\bottomrule
\end{tabularx}
\end{center}
Every check passed. The heterogeneous tests use eight seeded parameter
choices per preorder; they are not exhaustive over all multiplicity vectors.
Their seed is $20260919$. Each nonempty size class is exhausted without
random sampling of the relations. The source already reports larger finite
checks for some of its conjectures \cite[Sections~4 and~7]{AC26}.
Our run is an independent regression test of the formulas and their
implementation, not a claim to improve those enumeration bounds.

Run the checks from the package root with
\begin{verbatim}
python3 code/verify.py --max-n 4 --out data/verification.json
\end{verbatim}
The default run took about $19$ seconds in the execution environment used
for this manuscript. This timing is descriptive, not a benchmark guarantee.
The JSON records the measured time, all check totals, and exact data for the
worked examples. Repeated runs can differ in timing but should not differ
in any mathematical output. The \texttt{--max-n 5} option is available but
was not run for this report and can be substantially more expensive.

\subsection{Complexity and limitations of the checker}
The relation generator examines $2^{n(n-1)}$ reflexive relations before
filtering for transitivity. For a bound $B$ on the coordinate sum, the raw
candidate lattice vectors number $\binom{B+n}{n}$. These are deliberately
simple exhaustive algorithms, not efficient methods for large preorders.
The exact formulas often give faster counts after $P_{\tau^*}$ is available,
but the independent brute-force route is useful for detecting sign,
orientation, and indexing mistakes.

No floating-point geometric test, numerical root finder, or external
computer-algebra system is used in verification. A passing finite run is
supporting evidence; the general results rest on the proofs in
\Cref{sec:weighted,sec:reciprocity,sec:words,sec:matrix}.

\section{Proof audit, boundaries, and further questions}\label{sec:audit}

\subsection{Dependencies and delicate points}
The main dependencies can be summarized without assuming any of the
conjectures under investigation:
\[
 \begin{gathered}
 \text{Minkowski description + Postnikov}
       \ \Longrightarrow\ \text{polynomial }\FF_\tau;\\
 \text{interior translation + Ehrhart reciprocity}
       \ \Longrightarrow\ \text{multivariate reciprocity};\\
 \text{reciprocity + binomial sign identity}
       \ \Longrightarrow\ \text{falling formula}
       \ \Longrightarrow\ \text{word formula};\\
 \text{copied graph + Postnikov}
       \ \Longrightarrow\ \text{heterogeneous duality}
       \ \Longrightarrow\ \text{matrix duality}.
 \end{gathered}
\]
The classical theorems are stated explicitly and credited. Their proofs are
not claimed as new, nor is their full theory reproduced.

\paragraph{Dimension.}
Interior reciprocity is applied only to a full-dimensional lattice polytope
with strictly positive integer parameters. Degenerate boundary parameters
are reached by polynomial identity, not by changing the reciprocity sign
or invoking the relative interior of a different-dimensional polytope.

\paragraph{The all-ones shift.}
The amount removed from $x(I)$ is exactly $\card I$, not the number of
preorder classes in $I$. This is why equivalence classes of size greater
than one are covered by the same proof. The slack parameter must decrease
by one as well; forgetting this shift produces a false formula.

\paragraph{Signs.}
In deriving the falling formula, the outer sign is $(-1)^n$ and each
summand has the same inner sign because
$a_0+\sum_ea_e=n$. The signs cancel only after both are present. In
particular,
$\EE_\tau(m,0)=(-1)^n\EE_\tau(-m-1,-1)$, with the sign retained when
$n$ is odd.

\paragraph{Dual orientation.}
The neighborhood of a left vertex is an \emph{upper} set. Consequently the
associated draconian vectors satisfy upper-set bounds, and therefore lie
in the polytope of the \emph{dual} preorder. The non-self-dual example in
\Cref{eq:fork-matrices} detects a missing reversal.

\paragraph{Copies and the root coordinate.}
In the heterogeneous argument the total sum of a draconian vector is
$\sumv{\qvec}$, one less than the number of right vertices, not
$\sumv{\pvec}$. Its root coordinate is a uniquely determined slack.
Zero copies are handled explicitly in the equivalence between subset and
filter inequalities.

\paragraph{What computations certify.}
For each enumerated preorder the exact interpolation checks a whole double
polynomial, not just a few unrelated evaluations. No finite list, however,
proves the statement for all preorders. That quantifier is handled by the
geometric and graph arguments.

\subsection{Consequences that should not be overclaimed}
A nonnegative word expansion does not by itself imply real-rootedness of
$h^*$: sums of real-rooted polynomials need not be real-rooted without
additional structure. Our proof supplies neither a compatible family nor
an interlacing argument. It also does not construct a shelling of the
lattice-point poset.

The support enumerator
$\sum_{\avec\in P_\tau}t^{\#\supp(\avec)}$ is different from
$h^*(\QQ_\tau,t)$. The source's conjectures about the former must not be
inferred from our theorem about the latter. Finally, the $q$-zeta refinement
requires control of a weight statistic that is absent from the unweighted
interior and graph counts.

\subsection{Concrete directions left by these proofs}
The graph proof counts draconian vectors but does not produce a simple,
canonical map directly between the two array sets. One can investigate
whether a particular triangulation in Postnikov's duality yields a useful
explicit algorithm. Such a construction would be stronger than the
cardinality identity, even though the latter is now proved here.

The two binomial expansions suggest looking for a direct bijection between
the multiset model in \eqref{eq:rising} and the subset model in
\eqref{eq:falling}. A bijection preserving additional statistics could be
relevant to the refined zeta question. We do not assume that such a
statistic-preserving bijection exists.

The centered-parity identity also provides a smaller space of possible
coefficients for double Ehrhart computations. It can be exploited to reduce
interpolation costs, although the included checker intentionally does not
assume it while reconstructing the polynomial.

\section{Conclusion}
The selected open problem has a short mechanism: a strict ideal inequality
becomes a weak one after subtracting the all-ones vector and reducing the
slack by one. Ehrhart reciprocity along integer rays turns that mechanism
into the desired identity on the entire parameter space. The same method
works with an independent capacity at each element and gives a useful
falling-binomial expansion.

A second mechanism, independent copying on the two sides of a bipartite
preorder graph, proves the full matrix duality and a heterogeneous
strengthening. Together these arguments prove Conjectures~7.2, 4.2, and~4.4
as formulated in \cite{AC26}, and the ordinary specialization of
Conjecture~4.9. The proofs have been audited for degenerate parameters,
equivalence classes, signs, and reversal of the order, and are accompanied
by exact reproducible checks. Their external validation and their priority
relative to other work remain separate questions.

\appendix
\section{A reusable interior-shift principle}\label{app:general}
The shortest argument is not specific to preorders. The following statement
isolates its scope and also explains why integrality cannot simply be
omitted.

\begin{proposition}\label{prop:general-shift}
Let $A$ be an integer matrix with $n$ columns and no zero rows, and define
\[
 K(\bvec;s)=\{\xvec\in\R_{\geq0}^n:A\xvec\leq A\bvec+s\one\}.
\]
Suppose this family has a polynomial lattice enumerator $F(\bvec,s)$ for
nonnegative integer parameters, and that for all positive integer
parameters it is a bounded full-dimensional lattice polytope. Then
\[
 F(-\bvec,-s)=(-1)^n F(\bvec-\one,s-1)
\]
as a polynomial identity.
\end{proposition}
\begin{proof}
Full dimensionality and the absence of zero normals mean that interior
points satisfy precisely the strict versions of all the defining
inequalities. At integer points, strictness of the upper constraints is
$A\xvec\leq A\bvec+(s-1)\one$. Substituting
$\xvec=\yvec+\one$ yields
$A\yvec\leq A(\bvec-\one)+(s-1)\one$, and strict coordinate positivity
yields $\yvec\geq\zero$. Apply the same raywise reciprocity and polynomial
identity argument as in \Cref{thm:reciprocity}.
\end{proof}

For preorder polytopes, $A$ consists of the indicator rows of nonempty
ideals. The substantive additional input is not the translation itself but
the lattice-polytope and polynomiality properties, proved using the
Minkowski representation. For arbitrary systems of $0$--$1$ inequalities,
those properties may fail. Thus the interior-shift calculation alone does
not justify a general polynomial identity for every set system.

\section{Artifact inventory and source audit}\label{app:inventory}
The archive contains the editable \texttt{article.tex}, its compiled PDF,
a build script, the exact verification program, the recorded JSON data and
console log, a referee-oriented checklist, and a source-audit note.
No external font files or third-party article PDFs are redistributed.

The bibliography cites the public source version used to identify the
conjectures, Postnikov's theorem in its original paper, and Stanley's
classical reciprocity treatment. The source-audit note records the access
date and identifies the exact proposition and conjecture numbers relevant
to the proofs. The ordinary reciprocal identity and its word consequence
are attributed to the source as conjectures; the proofs and weighted
extensions in this manuscript are presented separately from those source
statements.

\begin{thebibliography}{9}

\bibitem{AC26}
Christos A. Athanasiadis and Fr\'ed\'eric Chapoton,
\emph{Polytopes and posets associated to preorders},
arXiv:2605.26916v1, May 26, 2026.
\url{https://arxiv.org/abs/2605.26916}.
Public version and PDF examined September 19, 2026.
Relevant items: Proposition~3.3, Theorem~2.1, Lemma~4.1,
Conjectures~4.2, 4.4, 4.9, Proposition~7.1, Conjecture~7.2,
and Remarks~4.10(a) and~7.3.

\bibitem{Postnikov09}
Alexander Postnikov,
\emph{Permutohedra, associahedra, and beyond},
International Mathematics Research Notices \textbf{2009} (2009), 1026--1106.
Preprint arXiv:math/0507163.
\url{https://arxiv.org/abs/math/0507163}.
Theorem~11.3 and Corollary~11.8.

\bibitem{Stanley74}
Richard P. Stanley,
\emph{Combinatorial reciprocity theorems},
Advances in Mathematics \textbf{14} (1974), 194--253.
\url{https://math.mit.edu/~rstan/pubs/pubfiles/22.pdf}.
Section~6: the Ehrhart--Macdonald law of reciprocity.

\end{thebibliography}
\end{document}
